\section*{Limitations}
\label{sec:limitations}

The scale of our experiments was constrained by the computational costs associated with multimodal sensing and multi-agent collaboration. In particular, processing heterogeneous sensor features through role-specified multi-agent collaboration incurs nontrivial inference overhead. To prioritize breadth of evaluation across diverse tasks, modalities, and baselines, we therefore conducted experiments on feasible subsets of each dataset rather than the full dataset. Future work should investigate the scalability of this framework to larger datasets and long-term sensing populations, as well as strategies for reducing inference overhead without sacrificing robustness.

This work is currently limited to classification tasks, as there is no established benchmark for evaluating LLM-based multimodal sensing across broader task types. We thus curated an evaluation suite by selecting datasets with objective ground-truth labels enable rigorous and reproducible accuracy-based evaluation. As a result, our evaluation does not cover human-centric reasoning or subjective judgment tasks. Extending the framework to such domains would require specialized data collection protocols and human annotations to assess the quality, coherence, and usefulness of generated reasoning.

This work focuses primarily on establishing modality-specific agents and the hybrid fusion protocol. As a result, we did not incorporate advanced prompting strategies, such as Self-Consistency~(SC) or iterative Self-Refinement~(SR), applied on top of \system{}. Similarly, while multi-agent collaboration protocols such as ReConcile could potentially improve performance by integrating confidence signals into the semantic fusion process, we prioritized isolating and validating the effectiveness of the core protocol without introducing additional components. These advanced combinations represent promising future directions for enhancing the reasoning capabilities of both modality agents and fusion agents.

% The hybrid fusion agent in \system{} delegates final arbitration between semantic and statistical fusion outcomes to an inference-time LLM meta-reasoner. This design avoids additional training, fine-tuning, or learned fusion parameters, thereby preserving a highly efficient single-round inference pipeline. Nevertheless, the arbitration process remains heuristic and may be sensitive to prompt formulation or distribution shifts, potentially limiting robustness under adversarial or highly noisy sensing conditions. Developing principled, training-free reliability estimates to further constrain meta-reasoning and improve robustness remains an important direction for future work.

We utilized non-fine-tuned LLMs to demonstrate the generalizability of our framework. We expect that further specialization, such as fine-tuning agents on modality-specific data or integrating Retrieval-Augmented Generation~(RAG) would further improve performance. By demonstrating that role specialization and structured data instructions alone yield significant gains, we establish a foundational step toward more complex multimodal systems. As sensor-specialized LLMs continue to emerge, \system{} can serve as a guideline protocol for designing effective collaboration among heterogeneous, modality-aware agents.

% The balancing mechanism between semantic and statistical fusion can be further optimized. Our current hybrid fusion strategy employs a single agent to reconcile these paths based on free-form reasoning. While this is the most direct implementation, future work could enhance this by measuring real-time signal quality or uncertainty. Providing the hybrid fusion agent with explicit metadata regarding the reliability of the input data would allow for a more precise and adaptive weighting of semantic reasoning versus statistical consensus.

\section*{Ethical Considerations}
\label{sec:ethical_considerations}

\noindent \textbf{Potential Risks.} \system{} is a general multimodal sensing framework that can be applied to domains including health-related tasks. In high-stakes applications such as clinical decision support or mental health assessment, incorrect predictions or misleading model-generated reasoning may lead to inappropriate user actions. As LLMs may produce hallucinated or poorly calibrated interpretations, deploying such systems without adequate safeguards and human oversight may introduce safety risks. We emphasize that \system{} is intended for research and exploratory use only. Any real-world deployment in safety-critical contexts should incorporate rigorous validation, regulatory compliance, and human-in-the-loop supervision. Further research is required to establish reliability guarantees, calibration mechanisms, and domain-specific safeguards before applying such systems to precision-sensitive applications.

\noindent \textbf{Use of LLMs.} We used LLMs for language polishing of the manuscript and code formatting.
